\chapter{Lab Solutions and Acceptance Guidance}
\index{lab solutions}\index{acceptance criteria}\index{capstone!solutions}%
This appendix provides reference solutions for the weekly Thursday labs and the six milestone
capstones. Each entry gives the approach, the critical steps, and the evidence that satisfies the
chapter's acceptance criteria. Use it to unblock yourself after a genuine attempt---not instead
of one: the labs are the curriculum, and the DP-600 measures exactly the judgment they build
\cite{microsoftDp600}.

\begin{notebox}
A solution counts only when the chapter's acceptance criteria are met with evidence: row counts,
screenshots, ledger rows, drill outputs. ``It ran once on my machine'' is the starting point of
each lab, never the end.
\end{notebox}

\section{Part I: Core Infrastructure (Weeks 1--4)}
\subsection{Week 1: Workspace, Lakehouse, and Shortcut}
\textbf{Approach.} Create the workspace on trial capacity, provision the lakehouse, and shortcut
to a public Blob container. \textbf{Critical evidence.} The shortcut path read from PySpark with
zero bytes uploaded; the screenshot shows the shortcut icon, not a copied folder.
\textbf{Common failure.} Confusing \texttt{Files/} uploads with shortcuts---if the data is
physically in OneLake, the lab's point was missed.

\subsection{Week 2: Medallion Transformations}
\textbf{Approach.} Bronze landing as-is, Silver typed and deduplicated, Gold conformed.
\textbf{Critical evidence.} Row-count reconciliation Bronze $\rightarrow$ Silver (accepted +
rejected = input), and Gold aggregates matching a direct source computation.
\textbf{Common failure.} Silent row drops between layers with no quarantine accounting.

\subsection{Week 3: Spark Notebook Transformations}
\textbf{Approach.} Parameterized notebook (\texttt{processing\_date}, \texttt{source\_system}),
explicit schemas, Delta writes. \textbf{Critical evidence.} The notebook reruns with a different
date parameter and no code edits. \textbf{Common failure.} Hardcoded personal workspace paths.

\subsection{Week 4: OPTIMIZE, V-Order, and Time Travel}
\textbf{Approach.} Baseline query timed, \texttt{OPTIMIZE ... VORDER}, accidental delete in a
training copy, \texttt{VERSION AS OF} recovery. \textbf{Critical evidence.} \texttt{DESCRIBE
HISTORY} showing optimize and delete; the recovered row count matching the pre-delete baseline.
\textbf{Common failure.} Running \texttt{VACUUM} before the recovery validation, destroying the
versions the lab exists to demonstrate.

\subsection{Milestone 1 Capstone}
The reference flow: Bronze shortcuts $\rightarrow$ parameterized Silver notebooks $\rightarrow$
Gold star schema $\rightarrow$ \texttt{OPTIMIZE VORDER}, with the operations note (retry unit,
retention, cadence) as the differentiating artifact. The acceptance criteria' decisive items are
the row-count reconciliation and the parameterization; reviewers check those first.

\section{Part II: Warehousing and BI (Weeks 5--8)}
\subsection{Week 5: COPY INTO Five Million Rows}
\textbf{Approach.} Staged Parquet under date-partitioned paths, \texttt{COPY INTO} per month,
manifest insert after each batch. \textbf{Critical evidence.} Five million rows, aggregates
matching the source profile, one manifest row per batch. \textbf{Common failure.} A malformed
OneLake DFS path---copy it from the lakehouse explorer (Appendix~B).

\subsection{Week 6: SSMS and the Silver-to-Gold View}
\textbf{Approach.} Entra-authenticated SSMS connection to the endpoint, the join view, the TVF,
the RLS objects, the deliberate (failing) write attempt. \textbf{Critical evidence.} Two-persona
row counts under \texttt{EXECUTE AS USER}. \textbf{Common failure.} Testing RLS as the
owner---where every policy passes trivially.

\subsection{Week 7: Relationship, Measure, and SCD2 Merge}
\textbf{Approach.} Surrogate-key one-to-many relationships, four DAX measures, the two-step
MERGE. \textbf{Critical evidence.} The changed customer shows exactly two rows, one current,
with contiguous validity dates; a second merge run changes nothing. \textbf{Common failure.}
Running the insert step before the expire step---order is correctness.

\subsection{Week 8: Direct Lake versus DirectQuery Benchmark}
\textbf{Approach.} Two identical models, identical visuals, measured timings, one deliberate
fallback. \textbf{Critical evidence.} The completed benchmark table and a captured fallback event
in the Capacity Metrics app. \textbf{Common failure.} Benchmarking before V-Order, then
concluding the wrong thing.

\subsection{Milestone 2 Capstone}
The reference flow: manifest-gated \texttt{COPY INTO} $\rightarrow$ stored-procedure star schema
(SCD2 customer) $\rightarrow$ custom semantic model with RLS $\rightarrow$ benchmarked Direct Lake
report. Reviewers check idempotence (re-run changes nothing), the failing DQ test, and RLS proven
\emph{as every role}, in that order.

\section{Part III: ETL Orchestration (Weeks 9--12)}
\subsection{Week 9: Metadata-Driven Pipeline}
\textbf{Approach.} Config table $\rightarrow$ Lookup $\rightarrow$ ForEach $\rightarrow$ contract
gate $\rightarrow$ Copy. \textbf{Critical evidence.} A renamed column rejected, an additive column
flowing, and a fourth table onboarded by \texttt{INSERT} alone. \textbf{Common failure.} ForEach
unbounded parallelism flattening the source.

\subsection{Week 10: Dataflow Gen2 with Quality Gate}
\textbf{Approach.} Named Power Query steps $\rightarrow$ staging destination $\rightarrow$
notebook assertions $\rightarrow$ quarantine $\rightarrow$ Gold. \textbf{Critical evidence.} The
deliberately broken payload quarantined with reasons; Gold untouched. \textbf{Common failure.}
The gate existing but not being wired into the pipeline---aspirational quality control.

\subsection{Week 11: Gateway Extraction}
\textbf{Approach.} Clustered gateway, least-privilege identity, watermark-incremental pipeline,
gateway-down drill. \textbf{Critical evidence.} The duplicate-free rerun after the drill.
\textbf{Common failure.} The gateway on a laptop that sleeps---the lab's own anti-pattern.

\subsection{Week 12: Master Orchestration Pipeline}
\textbf{Approach.} Dataflow $\rightarrow$ gate notebook $\rightarrow$ ledger, On Failure $\rightarrow$
alert with runbook link. \textbf{Critical evidence.} The deliberate-gate-failure alert firing with
the link, and one ledger row per outcome. \textbf{Common failure.} Wiring On Completion where On
Success was meant---partial failures publishing to Gold.

\subsection{Milestone 3 Capstone}
The reference flow: config-driven ingestion $\rightarrow$ dataflow cleansing $\rightarrow$ notebook
features $\rightarrow$ gated Gold, under one master pipeline with alerting. The decisive artifacts
are the deliberate-failure drill evidence and the watermark discipline (advance on success only).

\section{Part IV: Real-Time Intelligence (Weeks 13--16)}
\subsection{Week 13: Eventstream and Producer}
\textbf{Approach.} Custom app source, Python producer with unique \texttt{event\_id}s, KQL
destination. \textbf{Critical evidence.} The dedup query disagreeing with the raw count after the
replay test---proving duplicates existed and were handled. \textbf{Common failure.} The
connection string committed to source.

\subsection{Week 14: Rolling Average and Anomalies}
\textbf{Approach.} \texttt{make-series} $\rightarrow$ \texttt{series\_fir} $\rightarrow$
\texttt{series\_decompose\_anomalies}, plus the materialized view. \textbf{Critical evidence.} The
injected hot streak flagged, normal noise quiet, every query time-bounded. \textbf{Common
failure.} \texttt{default=0} silently inventing readings for offline devices.

\subsection{Week 15: Sustained-Threshold Alert}
\textbf{Approach.} Per-device rule, three-consecutive-minutes condition, cooldown, resolved
notification. \textbf{Critical evidence.} Two-minute breach fires nothing; sustained breach fires
exactly once. \textbf{Common failure.} The rule reading the raw stream, double-firing on replay.

\subsection{Week 16: Dashboard and Reverse ETL}
\textbf{Approach.} Bounded parameterized tiles at decision cadence; Dataverse upsert on the
alternate key. \textbf{Critical evidence.} The second reverse-ETL run changing zero record
counts. \textbf{Common failure.} Refresh cadence set to platform maximum rather than decision
cadence.

\subsection{Milestone 4 Capstone}
The reference flow: faker producer $\rightarrow$ Eventstream $\rightarrow$ deduped KQL $\rightarrow$
impossible-travel query $\rightarrow$ Reflex $\rightarrow$ Teams, with the replay drill producing
zero duplicate alerts. Reviewers check the dedup-before-rule ordering and the freshness indicator
on the dashboard.

\section{Part V: Data Science and ML (Weeks 17--20)}
\subsection{Week 17: Data Wrangler Export}
\textbf{Approach.} Visual cleaning $\rightarrow$ exported code $\rightarrow$ parameterized cell
$\rightarrow$ as-of-dated feature table. \textbf{Critical evidence.} Rebuilds for three historical
dates, identical on rerun. \textbf{Common failure.} The windowed feature computed past its as-of
date---leakage by construction.

\subsection{Week 18: Autologged Random Forest}
\textbf{Approach.} Six tracked runs, fixed snapshot, registry round-trip by alias.
\textbf{Critical evidence.} Every run tagged with dataset and commit; the model loading from
\texttt{models:/.../Staging} in a fresh notebook. \textbf{Common failure.} Comparing runs across
different data snapshots and calling the difference model quality.

\subsection{Week 19: SynapseML Sentiment}
\textbf{Approach.} Key-Vault-resolved credentials, batched calls, raw responses persisted, Gold
with \texttt{model\_version}. \textbf{Critical evidence.} The 50-row sample inspected before
scale-up; the measured cost per thousand rows. \textbf{Common failure.} A key literal in the
notebook---an incident once Week~23 syncs it.

\subsection{Week 20: Batch Scoring at 100k}
\textbf{Approach.} Alias-loaded UDF, partition-idempotent write, distribution gate.
\textbf{Critical evidence.} The injected skew caught by the gate; the rerun replacing rather than
duplicating. \textbf{Common failure.} Scoring against ad hoc features instead of the governed
feature table---training-serving skew by design.

\subsection{Milestone 5 Capstone}
The reference flow: leakage-safe features $\rightarrow$ autologged training $\rightarrow$ registry
$\rightarrow$ gated daily scoring $\rightarrow$ Direct Lake report, with drift triggering retraining
but never promotion. The decisive artifacts are the three failing-case tests and the
human-gated promotion evidence.

\section{Part VI: Governance and Exam Prep (Weeks 21--24)}
\subsection{Week 21: Roles and RLS}
\textbf{Approach.} Contributor for the builder, item-level for the analyst, RLS tested per
persona per engine path. \textbf{Critical evidence.} The access matrix with the Spark path's
result documented, not assumed. \textbf{Common failure.} SQL-layer security declared complete
while the Spark path reads everything.

\subsection{Week 22: Labels, Lineage, and Crypto-Shredding}
\textbf{Approach.} Label inheritance verified, lineage captured, both erasure strategies
demonstrated. \textbf{Critical evidence.} \texttt{DELETE} rows readable via time travel; shredded
subject unreadable everywhere including history. \textbf{Common failure.} Keys resolved from
anywhere but the vault.

\subsection{Week 23: Git and Deployment Pipeline}
\textbf{Approach.} Dev bound to Git, three-stage pipeline with rules, promotion through stages
only. \textbf{Critical evidence.} The rehydration path documented or executed; the secret scan
clean. \textbf{Common failure.} Deployment rules left unreviewed while silently rewiring
production references.

\subsection{Week 24: Mock Exams and Remediation}
\textbf{Approach.} Two timed, closed-book mocks scored by domain; misses classified; weak domains
remediated by rebuilding labs. \textbf{Critical evidence.} Two consecutive 80\%+ mocks with no
domain below 70\%. \textbf{Common failure.} Reviewing answer keys instead of rebuilding labs.

\subsection{Milestone 6 Capstone}
The reference flow: three workspaces, Git-bound Dev, staged pipeline with rules, PII masked before
Gold, labels inheriting, RLS tested, DR plan naming its three mechanisms. The decisive artifacts
are the clean-environment redeploy and the DR plan---the two things that cannot be faked in a
screenshot.

\begin{tipbox}
For portfolio use, each capstone README should lead with the problem statement and architecture
diagram, include clean-environment setup, screenshots of the pipeline run, dashboard, and alert,
the capacity cost estimate, and a lessons-learned section naming what broke. Interviewers probe
trade-offs: be ready to defend Lakehouse versus Warehouse, Import versus Direct Lake, and batch
versus streaming at ten times the scale.
\end{tipbox}
